% ID: SEL-D0017-BD-E07
% ESTADO: cotejada. Segunda pasada visual completa y borrador compilado revisados.
% Fuente: Selectividad/matematicas_aplicadas/matematicas_aplicadas_P_E.pdf
% SHA-256: bc88bfeb0f184eb0c806c4efe04e84c24f3e5fcef9caaa12e94fd48d9f67a389
% Página: 2; bloque D; ejercicio original 7.

El tiempo de desfase, en minutos, entre la hora de paso programada de un
autobús por cierta parada y la hora real a la que pasa, sigue una distribución
Normal de media desconocida y varianza 4. Se observa el paso del autobús
por la parada en 10 ocasiones elegidas al azar, registrándose los siguientes
desfases:
\[
4.7\quad 2.1\quad 3.6\quad 5.4\quad 0.0\quad 4.2\quad 4.0\quad
-0.2\quad 1.9\quad 5.2
\]
\begin{enumerate}[label=\alph*)]
  \item \textbf{(1.25 puntos)} Obtenga un intervalo de confianza al 97\%
  para el desfase medio en la hora de paso del autobús.
  \item \textbf{(1.25 puntos)} \textquestiondown{}Qué tamaño muestral mínimo
  sería necesario para estimar el desfase medio con un error inferior a
  30 segundos y un nivel de confianza del 95\%? \textquestiondown{}Cómo
  variaría dicho tamaño muestral si se aumentara el nivel de confianza?
\end{enumerate}
